\subsection{Optimization for Deep Learning}

The old optimization notes are migrated here only where the algorithms are used as training methods for neural networks.  General numerical optimization belongs to Part 18.

\begin{definition} \label{def:sl-neural-empirical-objective} \hlt{Neural Empirical Objective}\\
For data $(x_i,y_i)_{i=1}^n$, network $f_\theta:x\mapsto f(x;\theta)$, loss $L$, neural empirical objective is $J_n(\theta)=n^{-1}\sum\limits_{i=1}^nL(f_\theta(x_i),y_i)$.\\
The population objective is $J(\theta)=\mathbb E[L(f_\theta(X),Y)]$ under the data-generating law.
\end{definition}

\begin{breakablealgorithm}
\caption{Minibatch Backpropagation Step}
\begin{algorithmic}
\Require Minibatch $B$, parameters $\theta$, learning rate $\eta$
\State Compute predictions $f(x;\theta)$ for all $(x,y)\in B$
\State Compute minibatch loss $\widehat{L}_B(\theta)$
\State Run reverse-mode automatic differentiation to compute $g=\nabla_{\theta}\widehat{L}_B(\theta)$
\State Update parameters $\theta \leftarrow \theta-\eta g$
\end{algorithmic}
\end{breakablealgorithm}

\begin{remark} \label{rem:sl-deep-optimization-challenges} \hlt{Deep Optimization Challenges}\\
For empirical loss $J_n(\theta)$, deep optimization is shaped by:
\begin{enumerate}[label=(\roman*)]
\setlength{\itemsep}{0pt}
\item \hlt{geometry}: large $\kappa(\nabla^2J_n)$, saddles with $\nabla J_n\approx0$ and indefinite $\nabla^2J_n$, and plateaus with $\norm{\nabla J_n}\approx0$.
\item \hlt{gradient propagation}: large or small $\norm{\frac{\partial h^{(L)}}{\partial h^{(k)}}}$ causes exploding or vanishing gradients.
\item \hlt{stochasticity}: $g_k=\nabla J_n(\theta_k)+\xi_k$, with $\mathbb E[\xi_k\mid\theta_k]=0$ and $\operatorname{Var}(\xi_k\mid\theta_k)\ne0$.
\item \hlt{stabilization}: initialization, normalization, gradient clipping, skip connections, and learning-rate schedules control scale, conditioning, and update variance.
\end{enumerate}
These choices are part of the training specification, not post hoc implementation details.
\end{remark}

\begin{remark} \label{rem:sl-second-order-methods} \hlt{Second-Order Methods}\\
Let $J:\R^d\to\R$, $g_k=\nabla J(\theta_k)$, and $H_k=\nabla^2J(\theta_k)$.
\begin{enumerate}[label=(\roman*)]
\setlength{\itemsep}{0pt}
\item \hlt{Newton}: choose $p_k$ from $H_kp_k=-g_k$ and set $\theta_{k+1}=\theta_k+\eta_kp_k$.
\item \hlt{Conjugate gradient}: solve $H_kp=-g_k$ approximately using only products $H_kv$.
\item \hlt{BFGS/L-BFGS}: use $p_k=-B_kg_k$, where $B_k\approx H_k^{-1}$; L-BFGS stores few curvature pairs.
\end{enumerate}
Large networks make explicit $H_k$ storage or factorization prohibitive, so these methods usually act as preconditioners, curvature approximations, or analysis tools.
\end{remark}
